\section{A nonsmoothable threefold and its deformation space}
\label{sec05:counterexample}

The three separate nonvanishing conditions on an $A_2$ summand cannot be
replaced by nonvanishing of the vector in that summand.  The example below
also shows that changing the choices of local components does not repair
this weaker condition.

\begin{theorem}\label{thm05:counterexample}
There is a normal projective Calabi--Yau threefold $X$ with exactly thirty
ordinary double points and two anticanonical $\dP_6$ cone points such that
$K_S$ contains a relation nonzero at every singular point for each of the
two mixed profiles $S=LP,PL$, but $X$ has no smoothing.  A smooth projective
crepant resolution $Y_0$ has
\[
 (h^{1,1}(Y_0),h^{2,1}(Y_0))=(23,21).
\]
\end{theorem}

Here $L$ denotes the $A_1$ line and $P$ the $A_2$ plane, in the order of
the two cone points.  We prove the nonsmoothability assertion directly
from nodal necessity and the simultaneous partial resolution of the
$A_2$ family.  Its proof is independent of the sufficiency assertion in
Theorem~\ref{thm04:criterion}.

\subsection{The threefold and its relation matrix}

Let $P\subset\ZZ^4\otimes\RR$ be the convex hull of the following
vertices.  All indices in this section start at zero.
\begin{equation}\label{eq05:vertices}
\begin{gathered}
\begin{adjustbox}{max width=\textwidth,center}
\begin{tabular}{r@{\quad}c@{\qquad}r@{\quad}c@{\qquad}r@{\quad}c}
\toprule
$i$&$v_i$&$i$&$v_i$&$i$&$v_i$\\
\midrule
0&$(1,0,0,0)$&1&$(0,1,0,0)$&2&$(0,0,1,0)$\\
3&$(0,0,0,1)$&4&$(1,-1,1,-1)$&5&$(0,0,1,-1)$\\
6&$(1,-1,0,0)$&7&$(-1,1,0,0)$&8&$(0,0,-1,1)$\\
9&$(-1,1,-1,1)$&10&$(0,0,0,-1)$&11&$(0,0,-1,0)$\\
12&$(-1,1,-1,0)$&13&$(0,-1,1,0)$&14&$(-1,0,1,0)$\\
15&$(0,-1,0,1)$&16&$(-1,0,0,1)$&17&$(-1,0,1,-1)$\\
18&$(0,-1,1,-1)$&19&$(-1,0,-1,1)$&20&$(0,-1,-1,1)$\\
21&$(0,-1,-1,0)$&22&$(0,-1,0,-1)$&23&$(-1,0,0,-1)$\\
24&$(-1,0,-1,0)$&&&&\\
\bottomrule
\end{tabular}
\end{adjustbox}
\end{gathered}
\end{equation}
This is the fan polytope: the ambient fan consists of the cones over its
proper faces, and the anticanonical Newton polytope is $P^\circ$.
The polytope $P$ is reflexive.  It has $23$ facets, $71$ primitive edges,
and $69$ two-faces: $37$ unimodular triangles, $30$ unit parallelograms,
and two unit-edge hexagons.  The latter have cyclic orders
\begin{equation}\label{eq05:hexagons}
 F_0=(13,15,20,21,22,18),\qquad
 F_1=(14,16,19,24,23,17)
\end{equation}
and centres $(0,-1,0,0)$ and $(-1,0,0,0)$.  These two centres and the
origin are the only lattice points of $P$ other than its vertices.
Every parallelogram and hexagon has dual edge of lattice length one.
The complete list of parallelograms, in the cyclic orders used below, is
\begin{equation}\label{eq05:nodes}
\begin{gathered}
\begin{adjustbox}{max width=\textwidth,center}
\begin{tabular}{r@{\quad}c@{\qquad}r@{\quad}c@{\qquad}r@{\quad}c}
\toprule
$j$&node face&$j$&node face&$j$&node face\\
\midrule
0&$(0,1,5,4)$&1&$(0,1,9,8)$&2&$(0,1,12,11)$\\
3&$(0,3,15,6)$&4&$(0,6,20,8)$&5&$(0,6,21,11)$\\
6&$(1,2,14,7)$&7&$(1,5,17,7)$&8&$(2,3,15,13)$\\
9&$(2,3,16,14)$&10&$(2,5,17,14)$&11&$(3,8,20,15)$\\
12&$(4,5,17,18)$&13&$(4,6,15,13)$&14&$(4,6,21,22)$\\
15&$(5,10,23,17)$&16&$(7,9,16,14)$&17&$(7,12,23,17)$\\
18&$(8,9,12,11)$&19&$(8,9,19,20)$&20&$(8,11,21,20)$\\
21&$(9,12,24,19)$&22&$(10,11,21,22)$&23&$(11,12,24,21)$\\
24&$(13,14,16,15)$&25&$(13,14,17,18)$&26&$(15,16,19,20)$\\
27&$(17,18,22,23)$&28&$(19,20,21,24)$&29&$(21,22,23,24)$\\
\bottomrule
\end{tabular}
\end{adjustbox}
\end{gathered}
\end{equation}
All these lattice assertions follow by computing the supporting
hyperplanes and their face incidences from \eqref{eq05:vertices}; the
listed singular faces also specify the local lattice tests.  In
particular, each parallelogram has primitive three-dimensional cone
lattice, and each triangle in the central subdivision of a hexagon is
unimodular.

Let $X$ be a general anticanonical hypersurface in the projective
Gorenstein toric Fano fourfold defined by this fan.  Nondegeneracy and
the dual-edge lengths give one node for each face in
\eqref{eq05:nodes} and one anticanonical $\dP_6$ cone for each face in
\eqref{eq05:hexagons}.  There are no other singularities.  The hypersurface
avoids the torus fixed points, and
\[
 \omega_X\cong\cO_X,\qquad H^1(X,\cO_X)=H^2(X,\cO_X)=0.
\]
Insert the two centre rays and choose a projective crepant simplicial
refinement retaining the central subdivisions of the hexagons.  Its
restriction gives a smooth projective crepant resolution $f:Y_0\to X$;
the remaining possible ambient singularities lie over torus fixed
points and do not meet $Y_0$.  These are the standard hypersurface and
resolution constructions of \cite[\S\S4.1--4.2]{Batyrev94}.  This choice
selects a small resolution at each node; independent arbitrary choices
of all thirty diagonals are not needed.

Write $E_i\cong\operatorname{Bl}_3\PP^2$ for the exceptional divisor over
$F_i$.  In each of the cyclic orders \eqref{eq05:hexagons}, mark its
boundary curves by
\begin{equation}\label{eq05:boundary}
 (C_0,\ldots,C_5)=
 (e_1,h-e_1-e_2,e_2,h-e_2-e_3,e_3,h-e_1-e_3).
\end{equation}
The intersection form on $(h,e_1,e_2,e_3)$ is
$\operatorname{diag}(1,-1,-1,-1)$.  With the root basis fixed in
Section~\ref{sec:deformation},
\[
 \ell=-h+e_1+e_2+e_3,\qquad A=-e_1+e_3,\qquad B=e_2-e_3,
\]
the boundary intersection rows are
\begin{equation}\label{eq05:rootrows}
\begin{aligned}
 r_\ell&=(-1,1,-1,1,-1,1),\\
 r_A&=(1,-1,0,1,-1,0),\\
 r_B&=(0,1,-1,0,1,-1).
\end{aligned}
\end{equation}
All three classes are perpendicular to $K_{E_i}$.

Define a $36\times25$ matrix $M$ as follows.  Its first thirty rows are
indexed by \eqref{eq05:nodes}: a cyclic square $(i,j,k,l)$ gives entries
$(1,-1,1,-1)$ in those columns and zero elsewhere.  These rows fix signed
generators $\gamma_j$ of the corresponding nodal link homology; replacing
a small resolution reverses the relevant sign.  The last six rows are
$r_\ell,r_A,r_B$ on $F_0$, followed by $r_\ell,r_A,r_B$ on $F_1$, placed
in the six vertex columns in the orders \eqref{eq05:hexagons}.  Equations
\eqref{eq05:nodes}--\eqref{eq05:rootrows} specify every entry of $M$.
The two additional columns for the centre rays are zero, because the
root classes pair trivially with $K_{E_i}$.

The ambient divisor intersections detect precisely the homological
relations.  Indeed, $P$ has $28$ lattice points and no facet-interior
lattice points.  The only two-faces with interior points are the two
hexagons, and their dual edges have no interior lattice points.
Batyrev's divisor formula therefore has zero correction term and gives
\begin{equation}\label{eq05:h11}
 h^{1,1}(Y_0)=28-5=23.
\end{equation}
The same construction of the divisor classes shows that the $27$ ambient
boundary divisors span $H^2(Y_0,\QQ)$
\cite[Theorem~4.4.2 and Corollary~4.5.1]{Batyrev94}.  Thus, writing a
tuple of local classes as
\[
 w=(n_0,\ldots,n_{29},\lambda_0,a_0,b_0,\lambda_1,a_1,b_1),
\]
its total class
\begin{equation}\label{eq05:totalclass}
 \sum_{j=0}^{29}n_j\gamma_j+
 \sum_{i=0}^1(\lambda_i\ell_i+a_i A_i+b_i B_i)
\end{equation}
vanishes in $H_2(Y_0,\QQ)$ if and only if $wM=0$.
Consequently
\[
 K=\{w\in\QQ^{36}:wM=0\}
\]
is the full homological relation space.  Its restriction to a profile
$S$ is $K_S$.  Whenever these coefficients are used as tangent
coordinates, they are normalized by the fixed global volume form and
Lemma~\ref{lem:localhodge}; independent choices of raw local equation
parameters are not implicit in $M$.

\subsection{A nonzero relation and the obstruction to smooth fibres}

For the profile $LP$, take
\begin{equation}\label{eq05:weakwitness}
\begin{aligned}
 (\lambda_0,a_0,b_0,\lambda_1,a_1,b_1)&=(3,0,0,0,1,0),\\
 (n_0,\ldots,n_{29})={}&(-2,-1,2,2,-2,1,1,-2,2,2,\\
 &\hphantom{(}-3,6,-1,1,-2,-4,-3,2,2,-4,\\
 &\hphantom{(}5,-4,-4,6,3,3,1,1,-1,4).
\end{aligned}
\end{equation}
Multiplication by the explicitly defined matrix gives $wM=0$.
Every nodal coefficient is nonzero, and the two cone components are
$3\ell$ on $E_0$ and $A$ on $E_1$, both nonzero.  By
\eqref{eq05:totalclass}, this is a
homological relation, so the weaker nonvanishing condition holds.  An
explicit relation on $PL$ has cone coordinates $(0,1,0,3,0,0)$; its nodal
coordinates agree with \eqref{eq05:weakwitness} for $0\leq j\leq24$ and
are
\[
 (n_{25},n_{26},n_{27},n_{28},n_{29})=(-1,5,2,-2,7).
\]
It too annihilates $M$ and is nonzero at every singular point.

Let $D_i$ be the ambient divisor associated with $v_i$, restricted to
$Y_0$, and set
\begin{equation}\label{eq05:divisors}
\begin{aligned}
 H_L&=-\sum_{i\in\{0,1,2,4,5,6,7,13,14,17,18\}}D_i,\\
 H_M&=\sum_{i\in\{2,3,13,14,15,16\}}D_i,\\
 H_P&=\sum_{i\in\{0,1,2,3,4,5,6,7,13,14,15,16,17,18\}}D_i.
\end{aligned}
\end{equation}
Pairing \eqref{eq05:totalclass} with these three divisors gives, for
every $w\in K$,
\begin{equation}\label{eq05:identities}
\begin{aligned}
 n_{13}-a_0+b_0-a_1+b_1&=0,\\
 b_0+b_1&=0,\\
 n_{16}+\lambda_0+\lambda_1&=0.
\end{aligned}
\end{equation}
These equalities can be checked directly from the four-entry nodal rows
and the six-entry rows \eqref{eq05:rootrows}.  In particular they are
identities on the full relation space, before any local component is
selected.  They imply the following restrictions:
\begin{equation}\label{eq05:forced}
 \left.n_{13}\right|_{K_{LL}}=0,\qquad
 \left.b_1\right|_{K_{LP}}=0,\qquad
 \left.b_0\right|_{K_{PL}}=0,\qquad
 \left.n_{16}\right|_{K_{PP}}=0.
\end{equation}
For a plane component the interpretation of $b_i$ is intrinsic to the
marking:
\begin{equation}\label{eq05:threecoefficients}
 a_iA_i+b_iB_i=-a_i e_1+b_i e_2+(a_i-b_i)e_3.
\end{equation}
It is the coefficient of the second of the three nodal curves in the
simultaneous partial resolution.

\begin{proof}[Proof of the nonsmoothability assertion in
Theorem~\ref{thm05:counterexample}]
Suppose that $X$ has a convergent one-parameter smoothing.  At each cone
point the classifying map from the disc factors through one of the two
local components: in the domain $\CC\{t\}$, the equations
$\lambda_i a_i=\lambda_i b_i=0$ force either $\lambda_i=0$ or
$a_i=b_i=0$.  Hence the smoothing has one of the four profiles.

We use the following form of nodal necessity.  A smoothing and a chosen
ordinary double point produce a homological relation on a crepant
resolution whose coefficient at that node is nonzero; at the remaining
cone points its components belong to the selected branch subspaces.
To recall the argument, the local vanishing sphere has a nonzero period
of the Calabi--Yau volume form, and so is nonzero in the homology of a
smooth fibre.  A dual global three-cycle has a boundary tuple on the
links.  The boundary exact sequence makes its total class zero on the
resolution, and its component at the chosen node is nonzero.  The
boundary classes of the other local fillings give the stated branch
restrictions.  This is Friedman's nodal period argument
\cite{Friedman86}, with the cone boundary description of
\cite[Theorem~6.1]{PaperIV}; it concerns the given smoothing, without a
condition on its Kodaira--Spencer class.

For $LL$, apply this argument to node $13$.  The first identity in
\eqref{eq05:identities} forces its coefficient to be zero, a contradiction.
For $PP$, apply it to node $16$ and use the third identity.

For $LP$ or $PL$, apply the simultaneous modification of
Proposition~\ref{prop03:a2family} at the cone assigned to $P$.  More
explicitly, its local family consists of rank-at-most-one matrices with
diagonal entries $z+d_1,z+d_2,z+d_3$.  Remembering the column line and
writing the matrix as $\xi\eta^{\mathsf T}$ gives, on $\xi_k=1$,
\[
 \xi_i\eta_i-\xi_j\eta_j=d_i-d_j,
 \qquad \{i,j,k\}=\{1,2,3\}.
\]
The modification is proper and flat over the two-dimensional base.  Its
central fibre replaces the cone by three nodes with small-resolution
curves $e_1,e_2,e_3$.  It is an isomorphism over every noncentral base
point, since the zero matrix can occur only when all diagonal differences
vanish.  Pulling it back along the hypothetical smoothing and gluing to
the identity off the central cone point therefore gives a proper flat
family with the same smooth noncentral fibres.  Its central threefold
$Z$ has the thirty old nodes, these three new nodes, and the remaining
cone on its $L$ branch.  Crepancy and rationality give
$\omega_Z\cong\cO_Z$ and $H^1(Z,\cO_Z)=0$.

Apply nodal necessity to the new node represented by $e_2$.  The resulting
boundary relation on $Y_0$ has a nonzero coefficient at $e_2$.  However,
$H_M$ pairs trivially with every old nodal curve and with the remaining
$A_1$ root.  On the modified exceptional surface,
\[
 H_M|_{E_i}=e_1+(h-e_1-e_2)=h-e_2,
 \qquad (H_M\cdot e_1,H_M\cdot e_2,H_M\cdot e_3)=(0,1,0).
\]
Pairing the relation with $H_M$ forces its $e_2$ coefficient to vanish.
This contradiction excludes both mixed profiles.  The argument applies
also to smoothing arcs with arbitrarily high contact with the local
discriminant.  Together with \eqref{eq05:weakwitness}, it proves the
claimed failure of the weaker condition.  The Hodge numbers are
\eqref{eq05:h11} and Lemma~\ref{lem05:locallytrivial} below.
\end{proof}

\subsection{Tangent dimensions and the four selected deformation spaces}

\begin{lemma}\label{lem05:locallytrivial}
For this threefold,
\[
 h^1(X,\Theta_X)=h^{2,1}(Y_0)=21.
\]
\end{lemma}

\begin{proof}
The polar $P^\circ$ has $26$ lattice points and no facet-interior lattice
points.  Every edge of $P$ is primitive, so the correction term in
Batyrev's formula for $h^{2,1}$ vanishes.  Thus
\[
 h^{2,1}(Y_0)=26-5=21
\]
by \cite[Corollary~4.5.1]{Batyrev94}.  The counts may be obtained directly
from $\langle m,v_i\rangle\geq-1$ using \eqref{eq05:vertices}.

We verify the tangent-sheaf comparison needed for the first equality.
A vector field at an isolated cone point vanishes at the vertex, and
its local flow preserves the vertex ideal.  It therefore lifts to the
vertex blow-up $\operatorname{Tot}(K_E)$.  At a node the local flow lifts
to the blow-up and preserves each of the two ruling components of its
exceptional quadric: a one-parameter flow cannot interchange these
discrete choices.  It consequently lifts to the chosen small
resolution.  Conversely a vector field on the resolution descends on
the regular locus and extends across the point by normality, or
equivalently by reflexivity of the tangent sheaf.  Hence
\begin{equation}\label{eq05:pushforwardtheta}
 f_*\Theta_{Y_0}=\Theta_X.
\end{equation}
This is the equivariant-resolution comparison of
\cite[Definition~1.11 and Lemma~1.15]{FL22a} in these local models.

We also have $R^1f_*\Theta_{Y_0}=0$.  At a degree-six point use the
filtration by powers of the zero-section ideal.  The tangent bundle
restricts as $\Theta_E\oplus K_E$, so its graded terms are
\begin{equation}\label{eq05:gradedtheta}
 \Theta_E\otimes\cO_E(-kK_E),\qquad
 K_E\otimes\cO_E(-kK_E),\qquad k\geq0.
\end{equation}
The second term has vanishing $H^1$: use Serre duality and rationality
for $k=0$, rationality for $k=1$, and Kodaira vanishing for $k\geq2$.
For the first term, twist the toric Euler sequence
\[
 0\longrightarrow\cO_E^{\oplus4}
 \longrightarrow\bigoplus_{j=0}^5\cO_E(C_j)
 \longrightarrow\Theta_E\longrightarrow0.
\]
At $k=0$, the sequence
$0\to\cO_E\to\cO_E(C_j)\to\cO_{C_j}(-1)\to0$ gives
$H^1(E,\cO_E(C_j))=0$.  For $k\geq1$, the divisor $C_j-kK_E$ is nef:
its intersections with the six generators \eqref{eq05:boundary} of
the effective cone of curves are nonnegative, since $C_j^2=-1$ and
$(-K_E)\cdot C_l=1$ for every $l$.  Toric vanishing gives
$H^1(E,\cO_E(C_j-kK_E))=0$.  Finally
$H^2(E,\cO_E(-kK_E))=0$ for every $k\geq0$, so the twisted Euler
sequence proves the other $H^1$ vanishing in \eqref{eq05:gradedtheta}.

At a nodal small resolution the exceptional curve has normal bundle
$\cO_{\PP^1}(-1)^{\oplus2}$, and the corresponding graded tangent
terms are sums of $\cO_{\PP^1}(k+2)$ and
$\cO_{\PP^1}(k-1)$ for $k\geq0$.  Their $H^1$ groups also vanish.
Induction on the infinitesimal neighbourhoods of the exceptional sets
and the theorem on formal functions now give
$R^1f_*\Theta_{Y_0}=0$; these ideal filtrations are cofinal with the
pullbacks of the maximal-ideal filtrations at the singular points.
Leray, \eqref{eq05:pushforwardtheta}, and contraction with the crepant
volume form yield
\[
 H^1(X,\Theta_X)\cong H^1(Y_0,\Theta_{Y_0})
 \cong H^1(Y_0,\Omega^2_{Y_0}),
\]
which proves the assertion.
\end{proof}

Row reduction of the matrix specified above gives
\begin{equation}\label{eq05:fullranks}
 \rk M=20,\qquad \dim K=16,\qquad
 \rk M_{\mathrm{nodes}}=19.
\end{equation}
Here $M_{\mathrm{nodes}}$ consists of the first thirty rows.  Its relation
space has dimension $11$, so projection of $K$ to the six cone
coordinates has dimension $16-11=5$.  The second identity in
\eqref{eq05:identities} therefore describes its entire image:
\begin{equation}\label{eq05:coneimage}
 \operatorname{im}(K\longrightarrow\QQ^6)
 =\{(\lambda_0,a_0,b_0,\lambda_1,a_1,b_1):b_0+b_1=0\}.
\end{equation}

For a profile $S$, let $M_S$ consist of the nodal rows and the selected
root rows, and let $\rho_i:K_S\to R_{i,S_i}$ be its projection to the
$i$th cone summand.  The ranks and resulting dimensions are
\begin{table}[htbp]
\caption{Relation spaces and selected deformation spaces of $X$.}
\label{tab05:profiles}
\begin{adjustbox}{max width=\textwidth,center}
\begin{tabular}{c r r r c c r}
\toprule
$S$&rows of $M_S$&$\rk M_S$&$\dim K_S$&$\rk\rho_0$&$\rk\rho_1$
 &$\dim\Def^S(X)$\\
\midrule
$LL$&32&19&13&1&1&34\\
$LP$&33&20&13&1&1&34\\
$PL$&33&20&13&1&1&34\\
$PP$&34&20&14&2&2&35\\
\bottomrule
\end{tabular}
\end{adjustbox}
\end{table}
In particular both cone projections are nonzero on every profile.
Theorem~\ref{thm04:smoothbase} applies and proves that all four selected
deformation spaces are smooth, with the dimensions in the final column.
These dimensions use Lemma~\ref{lem05:locallytrivial} and the exact
sequence of Theorem~\ref{thm:D}.

For the unrestricted deformation space the same theorem gives
\begin{equation}\label{eq05:globaltangent}
 0\longrightarrow H^1(X,\Theta_X)
 \longrightarrow T_0\Def(X)
 \longrightarrow K\otimes_\QQ\CC\longrightarrow0.
\end{equation}
Thus $\dim T_0\Def(X)=37$.  The kernel of its projection to the two
cone tangent spaces has dimension $21+11=32$.  The second summand in
this count consists of global tangent directions whose local
components are supported at the old nodes.

\subsection{The unrestricted analytic deformation germ}

\begin{theorem}\label{thm05:germ}
For the threefold of Theorem~\ref{thm05:counterexample} there is an
isomorphism of convergent local rings
\begin{equation}\label{eq05:normalform}
 \cO_{\Def(X),0}\cong
 \frac{\CC\{z_1,\ldots,z_{32},x,y,u,v,b\}}
 {(xy,xb,uv,ub)}.
\end{equation}
The completed local ring has the same presentation with formal power
series.  The four selected deformation spaces are precisely its
irreducible components.  They are smooth, and their ideals in the
ambient smooth germ are
\begin{equation}\label{eq05:componentideals}
 \begin{aligned}
 P_{LL}&=(y,v,b),& P_{LP}&=(y,b,u),\\
 P_{PL}&=(x,v,b),& P_{PP}&=(x,u).
 \end{aligned}
\end{equation}
In particular the full deformation germ is reduced, although its
deformation functor is obstructed.
\end{theorem}

\begin{proof}
Choose coordinates on the two analytic local bases so that
\begin{equation}\label{eq05:localbases}
 \cO_{\Def(X,p_i),0}=
 \CC\{\lambda_i,a_i,b_i\}/(\lambda_i a_i,\lambda_i b_i)
\end{equation}
and their differentials are the Hodge-normalized coordinates used in
$M$.  This normalization is possible by an invertible linear change
on the line and on the plane separately; it preserves the displayed
ideal.  Fix the local classifying maps used to define $\Def^S(X)$.

By \eqref{eq05:coneimage} and \eqref{eq05:globaltangent}, the pulled-back
functions
\[
 x=\lambda_0,\quad y=a_0,\quad u=\lambda_1,\quad v=a_1,\quad b=b_0
\]
have independent differentials.  Complete them to a minimal analytic
embedding with $32$ further coordinates $z_1,\ldots,z_{32}$, and write
\[
 \cO_{\Def(X),0}=R/I,\qquad
 R=\CC\{z_1,\ldots,z_{32},x,y,u,v,b\},\qquad
 I\subseteq\mathfrak m^2,
\]
where $\mathfrak m$ is the maximal ideal of $R$.
Choose a representative $H\in R$ for the negative of the pullback of
$b_1$.  Equation~\eqref{eq05:coneimage} says exactly that
$H=b+O(\mathfrak m^2)$.  The local equations
\eqref{eq05:localbases} imply
\begin{equation}\label{eq05:localideal}
 F:=(xy,xb,uv,uH)\subseteq I.
\end{equation}
We will prove equality and then remove the higher terms of $H$ by a
convergent change of coordinates.

Let $\operatorname{in}(I)$ denote the ideal of lowest nonzero homogeneous
parts of elements of $I$ in
$\operatorname{gr}_{\mathfrak m}R=\CC[z_1,\ldots,z_{32},x,y,u,v,b]$.
The four elements in \eqref{eq05:localideal} give
\[
 J:=(xy,xb,uv,ub)\subseteq\operatorname{in}(I).
\]
On the other hand, each of the four selected germs is a smooth closed
subgerm of $\Def(X)$, by Table~\ref{tab05:profiles} and
Theorem~\ref{thm04:smoothbase}.  Its tangent space has the corresponding
ideal in \eqref{eq05:componentideals}; this follows from
\eqref{eq05:globaltangent} and the linear relation $b_1=-b_0$.
If $g\in I$, substitute a regular local parametrization of one of these
smooth germs into $g$.  Since the result vanishes, the lowest homogeneous
part of $g$ vanishes on its tangent space.  Therefore
\[
 \operatorname{in}(I)\subseteq
 P_{LL}\cap P_{LP}\cap P_{PL}\cap P_{PP}.
\]
The elementary monomial-ideal identity
\begin{equation}\label{eq05:primeintersection}
 (y,v,b)\cap(y,b,u)\cap(x,v,b)\cap(x,u)
 =(xy,xb,uv,ub)
\end{equation}
now proves
\begin{equation}\label{eq05:initialideal}
 \operatorname{in}(I)=J.
\end{equation}
For example, \eqref{eq05:primeintersection} can be checked monomial by
monomial: membership in all four coordinate primes forces divisibility
by at least one of $xy,xb,uv,ub$, and each of those four monomials
belongs to every prime.

Since $J$ is an intersection of prime ideals, the associated graded
ring of $R/I$ is reduced.  Its maximal-ideal filtration is separated by
Krull's intersection theorem.  A nonzero nilpotent would have a nonzero
initial form, and a power of that initial form would be zero, which is
impossible.  The same argument applies to the completion, whose
associated graded ring is unchanged.  Thus both local rings are
reduced.  This step uses the full equality \eqref{eq05:initialideal};
smoothness of the four selected germs by itself would not exclude an
additional nilpotent structure.

Equation~\eqref{eq05:initialideal} also proves $I=F$ in the convergent
ring.  Indeed, for $g\in I$, express its lowest homogeneous part as a
homogeneous polynomial combination of $xy,xb,uv,ub$.  Subtract the same
polynomial combination of $xy,xb,uv,uH$.  The remainder is in $I$ and
has strictly higher order.  For every integer $N$ a finite iteration
therefore gives $g\in F+\mathfrak m^N$.  Krull intersection in the
Noetherian local ring $R/F$ implies $g\in F$.  Hence
\begin{equation}\label{eq05:exactideal}
 I=(xy,xb,uv,uH).
\end{equation}
No convergence of an infinite sequence of correcting coefficients is
needed for this ideal equality.

It remains to straighten $H$.  The scheme-theoretic $LL$ restriction
imposes $y=b=v=H=0$.  By \eqref{eq05:exactideal}, its local ring is
\[
 R/(y,b,v,H)\cong
 \CC\{z_1,\ldots,z_{32},x,u\}/(H|_{y=b=v=0}).
\]
This selected germ has dimension $34$, equal to the dimension of the
convergent power-series ring on the right.  That ring is a domain, so
the displayed restricting function must vanish identically.  It follows
that
\begin{equation}\label{eq05:Hdecomposition}
 H=cb+py+qv,\qquad c(0)=1,\quad p(0)=q(0)=0,
\end{equation}
with convergent coefficient functions $c,p,q$.  In particular $c$ is a
unit.  Put $b'=cb+py=H-qv$.  Its linear part is $b$, so the analytic
inverse function theorem allows $b'$ to replace $b$ as an ambient
coordinate, with the other $36$ coordinates fixed.  The identities
\[
\begin{aligned}
 xb'&=c(xb)+p(xy),& ub'&=uH-q(uv),\\
 c(xb)&=xb'-p(xy),& uH&=ub'+q(uv)
\end{aligned}
\]
give
\[
 (xy,xb,uv,uH)=(xy,xb',uv,ub').
\]
Renaming $b'$ as $b$ proves \eqref{eq05:normalform} by a convergent
coordinate change.

The change also identifies the scheme-theoretic profile restrictions.
Their ideals before the change are
\[
 (y,b,v,H),\quad(y,b,u),\quad(x,v,H),\quad(x,u),
\]
in the order $LL,LP,PL,PP$.  Equation~\eqref{eq05:Hdecomposition} and the
invertibility of $c$ send these ideals to those in
\eqref{eq05:componentideals}.  The four coordinate primes are pairwise
incomparable, and \eqref{eq05:primeintersection} is their intersection.
They are therefore exactly the irreducible components.
Finally a tangent with $x=y=1$ and $u=v=b=0$ cannot extend to second
order, because $xy$ has a nonzero coefficient of $t^2$.  Thus the full
deformation functor is obstructed.
\end{proof}

The $32$-dimensional smooth factor in \eqref{eq05:normalform} is not the
locally trivial deformation space of $X$.  Its tangent space is the
kernel of localization at the two cone points; it contains the $21$
locally trivial directions and $11$ additional independent directions
with local components at the old nodes.  The final coordinate change
has identity linear part, so this tangent interpretation is preserved.

\begin{corollary}\label{cor05:intersections}
All scheme-theoretic intersections of the components in
Theorem~\ref{thm05:germ} are smooth.  The pairwise dimensions are
\begin{center}
\begin{adjustbox}{max width=\textwidth,center}
\begin{tabular}{c c c c c c c}
\toprule
pair&$LL,LP$&$LL,PL$&$LL,PP$&$LP,PL$&$LP,PP$&$PL,PP$\\
\midrule
dimension&33&33&32&32&33&33\\
\bottomrule
\end{tabular}
\end{adjustbox}
\end{center}
Every intersection of three or four distinct components is the same
smooth $32$-dimensional germ $x=y=u=v=b=0$.
\end{corollary}

\begin{proof}
Add the corresponding coordinate ideals in
\eqref{eq05:componentideals}.  Each sum is generated by a subset of the
five coordinates, giving the dimensions stated.  No reduction of these
intersection schemes is required.
\end{proof}

\begin{corollary}\label{cor05:integrability}
A global first-order deformation of $X$ is the tangent of a convergent
curve in $\Def(X)$ if and only if its five coordinates in
\eqref{eq05:normalform} satisfy
\[
 xy=xb=uv=ub=0.
\]
Equivalently, existence of a second-order extension of a tangent is
sufficient for existence of a convergent deformation with that tangent.
Each of the four components also contains a curve whose sufficiently
small noncentral fibres have exactly one ordinary double point and no
other singularities.
\end{corollary}

\begin{proof}
The quadratic equations give the necessity for a tangent by taking the
coefficients of $t^2$.  Conversely, a vector satisfying them belongs to
one of the four linear components in \eqref{eq05:componentideals}.
Its straight line in the convergent coordinates of
\eqref{eq05:normalform} integrates it.

For the last assertion, on each $K_S$ the coefficient in
\eqref{eq05:forced} is the only identically vanishing functional among
the old nodal coefficients, the selected line coefficients, and the
three separate plane coefficients \eqref{eq05:threecoefficients}.
This follows by row reduction of $M_S$ using the printed rows.
Avoiding the kernels of the other finitely many functionals gives a
tangent on the smooth selected germ with all these other coefficients
nonzero.  Integrate it there.  Every other local singularity then
smooths for sufficiently small nonzero parameter.  On $LL$ the remaining
point is old node $13$, and on $PP$ it is old node $16$.  On $LP$ or $PL$
the remaining plane parameter is noncentral; its only possible
singularity is the ordinary node on the corresponding discriminant
line, as in Proposition~\ref{prop03:a2family}.  Were the final point also
to smooth, the resulting curve would give a smoothing of $X$, contrary
to Theorem~\ref{thm05:counterexample}.  Compactness and openness of
smoothness exclude additional singularities away from the original
ones.
\end{proof}

The assertion about a prescribed tangent does not say that every
prescribed higher jet extends.  For example, $x=y=t^2$ and
$u=v=b=z_1=\cdots=z_{32}=0$ gives a deformation over
$\CC[t]/(t^4)$ which has no extension with those coefficients over
$\CC[t]/(t^5)$, since $xy=t^4$.  Its zero tangent nevertheless integrates.
Likewise, the four minimal equations in \eqref{eq05:normalform} do not
assert that $\operatorname{Ext}^2(\Omega_X^1,\cO_X)$ has dimension four.

The finite calculations used in this section are the face and lattice
counts for \eqref{eq05:vertices}, the divisor intersections
\eqref{eq05:identities}, and the ranks of the matrices specified by
\eqref{eq05:nodes}--\eqref{eq05:rootrows}.  Reproduction files and exact
rational outputs are supplied with the paper.\footnote{The geometric
data and intersection calculations are in
\texttt{a2\_branch\_selection\_counterexample.py}; the Hodge and selected
kernel calculations are in \texttt{branch\_smoothness\_inputs.py}; and
the unrestricted kernel and ideal calculations are in
\texttt{counterexample\_deformation\_germ.py}, all under
\texttt{certificates/}.  The independent standard-library calculation
\texttt{a2\_counterexample\_rational\_replay.py} reproduces the rational
matrix identities and ranks.}
The analytic ideal equality and coordinate change are proved above;
they do not follow from a finite truncation of the deformation equations.
